\section{Activation Grid and Protocol}

\subsection{Forecasting setup}

Given a multivariate input window \(X_{t-L+1:t}\in\mathbb{R}^{L\times C}\), the model predicts the next \(H\) target values \(Y_{t+1:t+H}\). We report mean squared error (MSE) as the primary metric and mean absolute error (MAE) as a secondary metric. For a configuration \(c\), dataset \(d\), horizon \(H\), and seed \(s\), let \(m(c,d,H,s)\) denote the MSE. Improvements over GELU are reported as paired relative reductions:
\begin{equation}
    \Delta(c,d,H,s)
    =
    \frac{m(\mathrm{GELU},d,H,s)-m(c,d,H,s)}
         {m(\mathrm{GELU},d,H,s)}.
\end{equation}
Positive values mean that configuration \(c\) improves over the same-seed GELU run under the same dataset-horizon protocol.

The main protocol follows the official-script style used for Informer ETT experiments: ProbSparse attention, time-feature embedding, \(d_{\mathrm{model}}=512\), \(d_{\mathrm{ff}}=2048\), eight heads, dropout \(0.05\), learning rate \(10^{-4}\), patience \(3\), batch size \(32\), and six training epochs. We call this an official-script-style protocol rather than a reproduction of every original Informer number, because our central comparison is internal: activation-only changes against our own matched GELU baseline.

\subsection{Activation families and placements}

Let \(z\) be the pre-activation tensor in an Informer feed-forward block. The hidden activation families are
\begin{equation}
    f(z) \in
    \{\mathrm{GELU}(z),\mathrm{ReLU}(z),\mathrm{Swish}(z),
      \tanh(z),\mathrm{softsign}(z),\tanh(z)+0.01\sin(z)\}.
\end{equation}
GELU is the baseline. ReLU and Swish are non-bounded controls. \texttt{tanh}, \texttt{softsign}, and \texttt{tanh\_sin001} test whether bounded sign-preserving nonlinearities are useful in long-horizon forecasting.

Placement is controlled independently from the activation family. The main full-hidden configurations replace the hidden activation in both encoder and decoder feed-forward blocks. We also test encoder-only \texttt{tanh}, decoder-only \texttt{tanh}, and hidden \texttt{tanh} combined with an output \texttt{tanh}. The output variant is treated separately because it changes the prediction range and can behave like target clipping rather than merely changing hidden representation.

\subsection{Datasets and horizon cells}

The Informer grid covers ETTh1, ETTh2, ETTm1, ETTm2, and eight Solar sites. ETTh cells use horizons \(24,48,168,336,720\); ETTm cells use \(24,48,96,288,672\); Solar cells use \(24,96\). This gives 36 dataset-horizon cells. Each cell contains nine activation configurations and three seeds, yielding 324 dataset-horizon-configuration cells and 972 matched seed-runs.

The sequence length, label length, encoder depth, decoder depth, and ProbSparse factor are horizon-specific and are recorded in the experiment manifest. Within a dataset-horizon cell, these fields are fixed across all activation configurations. This design is intentionally narrower than a general hyperparameter search: it asks whether the activation choice matters when everything else is matched.

\subsection{Aggregation and evidence policy}

The main tables use only rows with batch size \(32\), six epochs, finite metrics, complete three-seed coverage, and matching protocol fields. Repeated remote rows are deduplicated by dataset, horizon, seed, protocol fields, and activation signature before aggregation. Rows from earlier small-batch Lee-style diagnostics, mismatched configs, incomplete cells, and non-finite runs are excluded from the main evidence.

For each dataset-horizon cell, configurations are ranked by mean MSE over the three seeds. Static rankings average these within-cell ranks over the 36 cells. Cell-level gain intervals are computed by first averaging paired-seed improvements inside each cell and then summarizing the 36 cell means. We use these intervals as descriptive uncertainty checks, not as a claim that the 36 benchmark cells are an i.i.d. sample from a larger population.

\subsection{Regime diagnostics}

To probe why different activations win in different cells, we compute input-window factors including volatility, realized volatility, volatility shock, jump intensity, max absolute change, tail ratio, trend consistency, reversal rate, mean-reversion strength, and spectral summaries. Summary-level factor bins cover the full completed grid. Sample-level oracles are available only for cells with saved prediction arrays.

We distinguish diagnostics from deployment. Same-split factor oracles are optimistic upper bounds because they observe the completed static runs. Leave-one-seed-out or held-out evaluations are stricter because they ask whether a factor rule generalizes beyond the seed used to estimate it. The current leave-one-seed-out evidence is too weak and localized to support a deployable activation router, so the paper treats regime factors as explanation and experiment-design tools.
